\chapter{Topological Spaces}
\lecture{1}{Tue 21 Jan 2025 11:01}{Syllabus Day}
Textbook not required but I already bought it wooo also biweekly homework on blackboard always due at 11:59pm. We are encouraged to type homework (awewsomeeee) also why does this one guy in the class look like a literal reddit mod LOL also karim showed up 7m late. ``I cannot give everyone an A'' great its competitive now

``You've probably heard that a mug and a donut are the same thing'' I love this class. Making this notion precise and rigorous is difficult, and by the end of the semester we will be able to prove this notion.

\begin{definition}[Topological Space]\label{dfn:1.1}
	Let $X\neq \varnothing$ be a set. A topology on $X$ is a collection of subsets $\mathcal{T} \subseteq \mathscr{P} (X)$ such that
	\begin{enumerate}
		\item $\varnothing,X\in \mathcal{T} $;
		\item $\mathcal{T} $ is closed under \emph{finite} intersection;
		\item $\mathcal{T} $ is closed under union.
	\end{enumerate}
	We call the pair $(X,\mathcal{T} )$ a \emph{topological space}.
\end{definition}
\begin{definition}[Open Set]\label{dfn:1.2}
	Let $(X,\mathcal{T} )$ be a topological space. Any subset $A\subseteq X$ such that $A\in \mathcal{T} $ is called an \emph{open set}.
\end{definition}

Note that there is a finite restriction on (2), but not on (3). We will discuss this later. Definiton \ref{dfn:1.1} is extremely abstract and seemingly arbitrary, so consider the following examples to build some intuition.

The most basic example is that of the \emph{standard topology on $\mathbb{R}$}. Define $\mathcal{T} $ as the collection of all unions of open intervals of the form $(a,b)\subseteq \mathbb{R}$, together with the empty set. Trivially this definition satisfies property (1), since $(-\infty,\infty)=\mathbb{R}$ is open, and we define $\varnothing\in \mathcal{T} $. It also satisfies (3) by our definition. Showing (2) is slightly more difficult and requires $\varepsilon $-$\delta $ to make precise, and we will cover it later when we have this machinery.

Consider the open sets of the form
\[
	\left( \frac{-1}{2^n},\frac{1}{2^n} \right) ,\quad n\in\mathbb{N}
\]
in the standard topology on $\mathbb{R}$. However, if we take the intersection
\[
	\bigcap_{n\in\mathbb{N}} \left( -\frac{1}{2^n},\frac{1}{2^n} \right) =\left\{ 0 \right\}  
.\]
This is \emph{not an open set} in the standard topology on $\mathbb{R}$, as it is not an open interval and not empty. This is why the finite criterion on (2) is important.

Consider a more simple example. On any set $X$, the trivial topology $\mathcal{T} =\left\{ \varnothing, X \right\} $ is a topology. We trivially have criteria (1). Consider a finite collection $\left\{ A_n \right\}_{n=0}^k$. There are two cases. If $A_n=\varnothing$ for some $n$, then
\[
	\bigcap_{n=0}^{k} A_n=\varnothing
.\]
Otherwise, we have
\[
	\bigcap_{n=0}^{k} A_n=X
,\]
since we must then have $A_n=X\forall n$. Since $\varnothing,X\in \mathcal{T} $, we have (2). The same argument applies for (3); if $X\in \left\{ A_n \right\}_{n=0}^k $, we have the union is $X$; otherwise we have the union is empty. The trivial topology is also called the minimal topology, as it is the smallest possible topology since it must contain these subsets. 

The next example is the discrete topology. Let $\mathcal{T} =\mathscr{P} (X)$ for some set $X$. This very trivially satisfies definition \ref{dfn:1.1}, and we will explain later why we call it the discrete topology. This is also called the maximal topology, since it is the largest possible topology.

Let $X=\mathbb{R}$. We will now define what is called the Zariski topology, or sometimes the finite compliment topology. Let $\mathcal{T} $ be defined as all sets with either finite compliment in $\mathbb{R}$, or the empty set:
\[
	\mathcal{T} \coloneqq \left\{ Y\subseteq \mathbb{R}\mid Y=\varnothing \lor \left| \mathbb{R}\setminus Y \right| <\infty \right\}
.\]
We will show this is a topology. By definition, $\varnothing\in \mathcal{T} $, and we have $\mathbb{R}\setminus\mathbb{R}=\varnothing $ is finite, so $\mathbb{R}\in \mathcal{T} $. The second example is more interesting. Consider a finite collection $\left\{ A_n \right\}_{n=1}^k \subseteq \mathcal{T} $. We want to show
\[
	\bigcap_{i=1}^{k} A_i
\]
is open. By our definition, a set is open if and only if it's compliment in $\mathbb{R}$ is finite. DeMorgan's Law states that
\[
	\left( \bigcap_{i=1}^{k} A_i \right)^\complement = \bigcup_{i=1}^{k} A_i^\complement
.\]
``I assume you learned this in high school''. Note that the maximum cardinality of this union occurs when $A_i \cap A_j = \varnothing$ for all $i\neq j$, so we have
\[
	\left| \bigcup_{i=1}^{k} A_i^\complement \right| =\sum_{i=1}^{k} \left| A_i^\complement \right|
.\]
A finite sum of finite numbers is finite. More generally, a finite union of finite sets is finite. Now we show the third criterion holds. Suppose $\left\{ A_i \right\}_{i\in I}\subseteq \mathcal{T} $. By definition, $\left| A_i^\complement \right| <\infty$ for all $i\in I$. We wish to show
\[
	\left(\bigcup_{i\in I} A_i\right)^\complement < \infty
.\]
DeMorgan's Law also states that
\[
	\left( \bigcup_{i\in I} A_i  \right)^{c} = \bigcap_{i\in I} A_i^{c} 
.\]
Since the intersection of finite sets is finite, we have that $\cup \left\{ A_i \right\}_{i\in I}$ is open. The Zariski topology is the topology used in algebraic geometry.
\begin{definition}[Neighborhood]\label{dfn:1.3}
	Let $X$ be a topological space and let $x\in X$. An open set $U$ such that $x\in U$ is called a \emph{neighborhood} of $x$.
\end{definition}

The intuition behind a neighborhood is the collection of points near $x$. This may not be true in a general topology, but it is true in many topologies.

\begin{theorem}\label{thm:1.1}
	Let $X$ be a topological space. Then $A\subseteq X$ is open if and only if for all $x\in A$, there exists a neighborhood $U$ such that $x\in U\subseteq A$.
\end{theorem}
\begin{proof}
	Suppose $A$ is open and let $x\in A$. Then $U=A$ is a neighborhood of $x$. Now suppose for all $x\in A$, there exists an open neighborhood $U_x\subseteq A$ of $x$. Note that
	\[
		A\subseteq \bigcup_{x\in A} U_x 
	\]
	because
	\[
		y\in A \implies  \exists U_y\text{ s.t. }y\in U_y \implies  y\in \bigcup_{x\in A} U_x
	.\]
	However, since every $U_x\subseteq A$, we have also that
	\[
		\bigcup_{x\in A} U_x \subseteq A
	\]
	since
	\[
		y\in \bigcup_{x\in A} U_x \implies \exists U_x \text{ s.t. }y\in U_x \subseteq A 
	.\]
	Therefore, $A$ is open by criteria (3) in definition \ref{dfn:1.1}.
\end{proof}
